\section*{Chapter 1. Theoretical foundations of fine-tuning \\ acceleration through importance analysis}
\addcontentsline{toc}{section}{Chapter 1. Theoretical foundations of fine-tuning \\ acceleration through importance analysis}
\setcounter{section}{1}

    The goal of the present study is to develop a method for accelerating the fine-tuning of Large Language Models. However, the development of a novel approach requires a comprehensive analysis of existing techniques and promising ideas in this field, which can be leveraged in future work. This chapter is dedicated to such a thorough analysis and review.
    
    \subsection{Definition of Transformer-based architecture}
    Let $\mathcal{M}_\Phi$ represent an autoregressive Transformer model with parameters $\Phi \in \mathbb{R}^n$. The model $\mathcal{M}_\Phi$ is constructed as a sequence of Transformer blocks, denoted by $\tau$. Each layer $\tau_i$ is comprised of two main components: a Multi-Head Attention module (MHA) (\ref{eq:mha}) and a Feed-Forward Neural Network (FFNN) (\ref{eq:mlp}). For an input $x \in \mathbb{R}^n$, the MHA component utilizes $h$ parallel Attention heads to process the sequence.

    \begin{equation}\label{eq:mha}
    \begin{aligned}
        head_i & = Softmax(\frac{xW_q (xW_k)^T}{\sqrt{d_h}})xW_v \\
        MHA (X) & = Concat (head_1, \dots , head_h)W_o, 
    \end{aligned}
    \end{equation}
    where $i \in \mathbb{R}$, $W_o \in \mathbb{R}^{d \times d}$ is an output projection, and \(W_q, W_k, W_v \in \mathbb{R}^{d \times d_h}\) are query, key, and value projections of the $i$-th head; $d_h$ is typically set to $d/h$. 
    
    FFNN module consists of two linear transformations with an activation function (e.g. ReLU) (\ref{eq:mlp}). 

    \begin{equation}\label{eq:mlp}
        FFNN(x) = ReLU(xW_{up} + b_{up})W_{down} + b_{down}, 
    \end{equation}
    where $W_{up} \in \mathbb{R}^{d \times d}$ is the up-projection matrix and $W_{down} \in \mathbb{R}^{d \times d}$ is the down-projection matrix. 

    Within the \llm{} there exist a concept of the  \textit{importance} of certain structures. In this context, \textit{importance} facilitates the identification of structures that can be removed without affecting model performance. Thus, understanding the \textit{importance} of a structure helps assess its impact on the model's overall quality.

    Fine-tuning is the process of turning a general-purpose model into a specialized one while updating all pre-trained model parameters over a downstream task dataset. During the fine-tuning process \llm{} parameters are updated untill the loss function reaches approximately the optimum, indicating, that the model has been converged \cite{allenzhu2019convergencetheorydeeplearning} and is ready for deployment. 

    Assume that the data intended for fine-tuning is structured as a dataset $D = \{(x_i, y_i), i = 1, \ldots, N\}$, where $x_i$ represents an input sample and $y_i$ is the corresponding label; each $x_i$ and $y_i$ consists of sequences of tokens. The response produced by the model for $x_i$ is expressed as $\mathcal{M}_\Phi(x_i)$. It is presumed that the model has already been pre-trained, so its parameters are initially set to $\Phi_0$.

    Full fine-tuning involves altering these pre-trained parameters to $\Phi_0 + \Delta \Phi$, aiming to optimize the fine-tuning objective for $\mathcal{M}_\Phi$ with respect to $\Phi$, as given by (\ref{eq:full_fine_tuning}):
    \begin{equation}
        \label{eq:full_fine_tuning}
        \max_{\Phi} \sum_{(x, y) \in D} \sum_{t = 1}^{|y|} \log\left(\mathcal{M}_\Phi(y_t \mid x, y_{<t})\right),
    \end{equation}
    where $|y|$ denotes the number of tokens in the label sequence, and the dimensionality of $\Delta \Phi$ is the same as that of the initial parameter vector $\Phi_0$ \cite{lora}.
    
    The resources required for full fine-tuning of a model are comparable to those needed for pre-training an \llm{}; the only difference is that the full fine-tuning necessi-tates a smaller amount of a domain-specific data. Therefore, the development of approaches, that allow to reduce the number of required resources is of crucial importance. 
\subsection{Key directions in Large Language Models fine-tuning \\ acceleration}
    Several modern techniques have been developed to minimize the computational and memory demands associated with full fine-tuning. These approaches include reducing the training data volume (data reduction) \cite{aloe}, reducing the data precision (quantization) \cite{qlora}, reducing the number of parameters in an \llm{} (pruning) \cite{model_pruning}, eliminating resource-intensive backpropagation (backpropagation elimination) \cite{mezo}, \cite{hinton2022forwardforwardalgorithmpreliminaryinvestigations}, and leveraging Parameter-Efficient Fine-Tuning methods \cite{lora}.

    \subsubsection{Data reduction}
        Data reduction refers to the practice of removing less influential examples from a training set while aiming to preserve the \llm{} accuracy. This concept is generally realized through three main strategies: dataset distillation or condensation, coreset selection, and importance sampling, each offering a unique way to compress the training dataset.

        In dataset distillation, synthetic examples are generated based on statistical characteristics extracted from the original data \cite{cazenavette2022datasetdistillationmatchingtraining}. This strategy seeks to replace the full dataset with a much smaller, yet still representative, set of artificial samples that embody the key data patterns. Coreset selection, on the other hand, involves heuristics-based identification and selection of a crucial subset from the initial dataset \cite{shim2021coresetsamplingefficientneural}, ensuring that essential samples are retained and training efficiency is enhanced. Importance sampling deals with prioritizing those samples which are expected to have the greatest impact on model learning, primarily by focusing on instances that cause significant parameter updates during training \cite{jiang2019acceleratingdeeplearningfocusing}. For instance, ALOE \cite{aloe} prunes data points according to their loss values, evaluated by a trained or pre-trained model, achieving on average up to 33\% faster fine-tuning with a quality drop of 1.5\% compared to \lora{}. However, these data reduction strategies introduce additional computational cost, as they require a complete forward pass through the dataset to estimate the importance of each sample \cite{NEURIPS2018_967990de}.  

    \subsubsection{Quantization}
        Quantization refers to the process of converting data with high numerical precision into a format with lower precision. Although this method is primarily employed to speed up inference tasks, recent research has investigated its potential to also make fine-tuning processes more efficient.

        A prominent example is QLoRA \cite{qlora}, which leverages a 4-bit NormalFloat format to store the static, pre-trained language model. Meanwhile, essential computational tasks -- such as matrix multiplications involving \lora{} adapters -- are performed using 16-bit BrainFloat precision. While this strategy significantly reduces memory requirements, it tends to be slower than the traditional \lora{} approach, owing to the overhead of quantizing and dequantizing \llm{} parameters. Furthermore, applying extremely low-bit quantization may introduce substantial quantization errors, which can negatively affect the initial phase of \lora{} fine-tuning \cite{han2024parameterefficientfinetuninglargemodels}. Another technique, PEQA \cite{peqa}, adjusts the quantization scales of quantized \llm{} while keeping the integer matrices intact. Although this can cut memory usage by up to 92\%, it often comes at the cost of a notable drop in model quality, around 5\% compared to \lora{}. AlphaTuning is yet another method that involves quantizing a pre-trained \llm{} and perform selective fine-tuning of certain quantized parameters for targeted tasks \cite{AlphaTuning}. This approach achieves a 38\% acceleration in fine-tuning and reduces memory consumption by 80\%, but with a minimum quality reduction of 5\%.
    \subsubsection{Pruning}
        Model pruning is a technique aimed at reducing the model by eliminating redund-ant parameters or substructures. This process is primarily categorized into three approaches: unstructured, semi-structured, and structured pruning \cite{model_pruning}.

        Unstructured pruning removes individual parameters based on importance crite-rion, achieving high theoretical sparsity and minimal quality degradation. SOTA methods in this area reports less than 5\% quality drop at 50\% unstructured sparsity \cite{xu2024besapruninglargelanguage}. However, its irregular memory access complicates hardware acceleration \cite{MLSYS2023_a10deb4d} and performance often matches or falls bellow dense counterparts due to overheads in parameter managements \cite{bambhaniya2024progressivegradientflowrobust}.

        Structured pruning targets entire model components (e.g. channels, layers, Trans-former blocks), making it hardware-friendly \cite{bambhaniya2024progressivegradientflowrobust}, but causing significant quality degradation due to removal of large parameter groups \cite{cheng2024surveydeepneuralnetwork}. SOTA methods in this area provide up to 13\% quality degradation at 50\% structured weight sparsity \cite{llm_pruner}. 

        Semi-structured pruning removes N out of every M consecutive weights and thus offers a favorable compromise between accuracy and hardware efficiecncy across various GPU architectures \cite{huang2024pruninglargelanguagemodels}. However, SOTA approaches in this field still provides up to 5\% quality degradation under 2:4 sparsification pattern \cite{yu2026efficientposttrainingpruninglarge}. 
        
        Moreover, with weight pruning the degradation in model performance extremely increases as the proportion of pruned parameters grows \cite{alanova2026motivatingnextgenacceleratorsflexible}. In addition, pruning process often introduces computational overhead on additional inference on calibration data or further training to recover the \llm{} performance \cite{sun2026maskprolinearspaceprobabilisticlearning}. 

    \subsubsection{Backpropagation elimination}
        Removing backpropagation entails substituting the conventional backpropagation algorithm with alternative learning approaches that are more computationally efficient. The main goal of this replacement is to substantially lower both memory consumption and processing time. Research in this area generally splits into two principal branches: zeroth-order techniques and methods based on local learning rules.

        Zeroth-order approaches estimate the gradient by performing only forward computations. An important approach within this category is the Memory-efficient Zeroth-order Optimizer (MeZO) \cite{mezo}, which offers an effective implementation of the ZO-SGD algorithm \cite{spall}. Research have demonstrated that MeZO can match the effectiveness of standard backpropagation fine-tuning on a range of tasks. Remarkably, MeZO is capable of cutting memory usage by up to twelve times and shortening training durations by as much as half. Nevertheless, it typically demands more training steps than traditional fine-tuning methods and only achieves satisfactory accuracy when used alongside a specifically designed prompt \cite{mezo}.

        The other main direction focuses on local learning mechanisms, especially the Hebbian learning rule. Inspired by the way the human brain is structured and operates, the Hebbian rule \cite{hebb} posits that synaptic connections between neurons become stronger when the neurons are activated together, and weaker when they are not. An efficient combination of localized learning rule with backpropa-gation provides an acceleration by up to 77\%, however introduces a substantial quality degradation of 3.04\% on average \cite{hebb_pricai}.  

    \subsubsection{Parameter-Efficient Fine-Tuning}

        Parameter-Efficient Fine-Tuning (\peft{}) has become a prominent strategy for accelerating the fine-tuning of \llms{}. This technique involves modifying only a limited portion of the model’s parameters, which drastically lowers both computational and memory requirements while maintaining, or sometimes even improving, predictive performance \cite{han2024parameterefficientfinetuninglargemodels}. In the context of \peft{}, $\Delta \Phi$ is encoded by a much smaller-sized set of parameters $\Theta$, so that $|\Theta| \ll |\Phi_0|$, $\Delta \Phi = \Delta \Phi (\Theta)$. Thus, \peft{} aims at maximizing the fine-tuning objective over $\Theta$ (\ref{eq:peft}). 

        \begin{equation}\label{eq:peft}
            max_{\Theta} \sum_{(x, y) \in D} \sum_{t=1}^{|y|} log (\mathcal{M}_{{\Phi}_0 + \Delta \Phi (\Theta)} (y_t | x, y_{< t}))
        \end{equation}

        A widely adopted method in industry is Low-Rank Adaptation (\lora{}) \cite{lora}. Rather than updating all the weights in the pre-trained model, \lora{} keeps the original model parameters fixed and introduces additional trainable low-rank matrices into each Transformer layer. This approach substantially cuts down the number of parameters that need to be adjusted during the fine-tuning. If we derive $W_0 \in \mathbb{R}^{d_1 \times d_2}$ as pre-trained weights of an arbitrary layer, $x$ as an input data, then, inference of this layer, denoted as $h$ can be expressed as \(h = W_0x\). \lora{} modifies the forward pass of the layer as follows (\ref{eq:lora}). 

        \begin{equation}\label{eq:lora}
            h = W_0x + \Delta W_x = W_0x + BAx,
        \end{equation}
        where $W_0, \Delta W \in \mathbb{R}^{d_1 \times d_2}, A \in \mathbb{R}^{r \times d_2}, B \in \mathbb{R}^{d_1 \times r}, r \ll (d_1, d_2)$. $A, B$ matrices are called \lora{} adapters. 

        By reducing the number of parameters involved in fine-tuning, \lora{} achieves up to 50\% savings in memory usage and can speed up the fine-tuning process by up to 150\% compared to traditional full fine-tuning. Despite these advantages, \lora{} has several limitations. Firstly, it does not allow for fine-grained modifications to the weights during fine-tuning \cite{dora}. Secondly, the initialization scheme for \lora{}’s adapters may slow down convergence \cite{pissa}. Thirdly, it treats all weight matrices equally, which ignores their varying importance \cite{adalora}. Lastly, \lora{} cannot fully match the performance of full fine-tuning, since its uniform updates across all layers prevents it from making larger adjustments where they matter most in the model \cite{Demidovskij2024}.
        
        To address these limitations, DoRA \cite{dora}, PiSSA \cite{pissa}, and AdaLoRA \cite{adalora} were further developed as promising improvements of \lora{}.
        
        \dora{} \cite{dora} approach splits the pre-trained weights $W_0$ into two distinct parts: \textit{magnitude} and \textit{direction} (\ref{eq:dora}). The \textit{magnitude} is a trainable vector $M = \|W_0\|_c$, representing the norm of each column of the weight matrix, while the \textit{direction} is defined as the normalized weight matrix. The direction component is fine-tuned using the \lora{} mechanism \cite{lora}.
        
        \begin{equation}\label{eq:dora}
            W = M \frac{W_0 + BA}{\|W_0 + BA\|_c}
        \end{equation}

        \dora{} provides up to 3\% of quality improvement on the fine-tuning of \llamaTwoSmall{} and \llamaTwoBig{} compared with \lora{}.

        \adalora{} \cite{adalora} approach assigns the parameter budget dynamically among the model’s weight matrices based on their estimated importance. It reformulates the incremental updates through Singular Value Decomposition (SVD), revising \lora{}’s forward computation as described in (\ref{eq:adalora}). 

        \begin{equation}\label{eq:adalora}
            h = W_0x + \Delta W_x = W_0x + P \Lambda Q,
        \end{equation}
        where $P \in \mathbb{R}^{d_1 \times r}$ and $Q \in \mathbb{R}^{r \times d_2}$ represent singular vectors of $\Delta W$, the diagonal matrix $\Lambda \in \mathbb{R}^{r \times r}$ contains the singular values $\{\lambda_i\}_{1 \leq i \leq r}$. 
        
        \adalora{} provides 1\% of quality improvement on DeBERTa-V3-base fine-tuning compared with \lora{}. However, both \dora{} and \adalora{} still incurs uniform updates of all \llm{} layers with a little attention to the critical initial projections \cite{Demidovskij2024}. 

        The \pissa{} approach \cite{pissa} enhances \lora{} by initializing its adapter matrices $A$ and $B$ with the leading principal components of the original weights W, obtained using SVD (\ref{eq:pissa}):

        \begin{equation}\label{eq:pissa}
            \begin{aligned}
                W & = USV^T, \\
                W^{res} & = U_{[:, r:]}S_{[r:, r:]}V^{T}_{[:, r:]}, \\
                A & = U_{[:, :r]}S^{\frac{1}{2}}_{[:r, :r]}, \\
                B & = S^{\frac{1}{2}}_{[:r, :r]}V^T_{[:, :r]},
            \end{aligned}
        \end{equation}
        where \(S\) is a diagonal matric, which represents the singular values in descending order. 
        
        This initialization ensures that the trainable adapters focus on the most informative directions of \(W\), thus closely mimicking the behavior of full fine-tuning of an \llm{} \cite{Demidovskij2024}. This method can yield up to a 21.05\% improvement in fine-tuning quality on \llamaTwoSmall{} compared to \lora{}.

    \subsubsection{Summary}
        Therefore, multiple methods have been developed to decrease the resource requirements of full fine-tuning.
        \begin{itemize}
            \item \textbf{Data reduction} aproaches accelerate training by minimizing the amount of data used, which decreases the computational effort required. However, these methods introduce extra overhead on importance estimation of samples  and lead to quality degradation due to removal of training samples from the dataset.
            \item \textbf{Quantization}  approaches reduce both memory consumption and training duration, but typically at the cost of a noticeable decline in model accuracy.
            \item  \textbf{Pruning} techniques are capable of greatly improving fine-tuning speed and lowering memory usage. However, these approaches may degrade model accuracy, introduce extra time for calibration or retraining, and often require specialized hardware to reach optimal efficiency.
            \item \textbf{Backpropagation elimination} methods significantly reduce computational de-mands by lowering memory and training time, but often require more iterations to achieve satisfactory accuracy or introduce quality degradation compared to backpropagation-only approaches.
        \end{itemize}

        Consequently, Parameter-Efficient Fine-Tuning has emerged as the most prominent solution. De facto standard method in this field, \lora{}, can speed up training by up to 150\% and cut memory requirements in a half, while maintaining performance comparable to full fine-tuning. Further efficiency gains are possible by strategically attaching adapters only to critical Transformer blocks within the \llm{}.

\subsection{Existing approaches for assessing the importance  within \\ Large Language Model}
    Assessing the importance within \llms{} can be applied at various granularities, from single parameter through Transformer blocks, up to the entire architecture.
    \subsubsection{Importance of separate weights}
        The Lottery Ticket Hypothesis \cite{frankle2019lotterytickethypothesisfinding} proposes the existence of a compact subnetwork -- often called the winning ticket -- embedded within a randomly initialized, overparameterized model. This subnetwork, when trained independently, can reach accuracy levels similar to the original network, sometimes even with equal or fewer training epochs. This phenomenon has been observed across different model architectures in both Computer Vision and Natural Language Processing directions. For example, a LeNet subnetwork retaining just 51.3\% of the original weights can often outperform the full model in fewer epochs. Similarly, in BERT-base, a highly sparse winning ticket (with 90\% of weights pruned) can still perform comparably to the full model within a standard deviation \cite{chen2020lotterytickethypothesispretrained}.

        Building on the idea of individual elements importance, the concept of systematic outliers in \llms{} has been introduced. These are values that substantially depart from the typical distribution \cite{an2025systematicoutlierslargelanguage}. Such outliers tend to bias model predictions towards frequent tokens, thereby affecting overall performance. Their emergence is most notable in the Multi-Head Attention module, where the SoftMax operation can amplify dominant attention weights exponentially for certain keys and values (\ref{eq:dominant_keys}), while suppressing others to near-zero (\ref{eq:other_keys}). This uneven attention distribution causes substantial updates in the $K$ and $V$ projection weights during backpropagation. When the Multi-Head Attention  outputs extreme values, these are further accentuated by non-linearities in the FFNN block, especially within the down-projection matrix. Thus, steep gradients and non-linear operations work together to foster the rise of outliers in activations and weights.
        \begin{equation}\label{eq:dominant_keys}
        A_{xj} = \frac{\exp(M)}{\exp(M) + (n - 1)} \approx 1
        \end{equation}
        \begin{equation}\label{eq:other_keys}
        A_{xk} = \frac{1}{\exp(M) + (n - 1)} \approx 0,
        \end{equation}
        let \(M = \frac{W_q (W_k^T)_j}{\sqrt{d_h}}, \frac{W_q (W_k^T)_k}{\sqrt{d_h}} = 0, \text{then} \frac{W_q (W_k^T)_j}{\sqrt{d_h}} \gg \frac{W_q (W_k^T)_k}{\sqrt{d_h}}\), $n$ is the number of elements.

        Several distinct categories of outliers have been identified \cite{an2025systematicoutlierslargelanguage}, including:
        \begin{enumerate}
            \item \textbf{Acitvation outliers:} Consider the output of a layer $h \in \mathbb{R}^{B \times H}$, where $B$ stands for the batch size and $H$ is the hidden dimension. The mean absolute activation value is defined as $\mu_h = \frac{1}{B \cdot H} \sum_{i, j} |h_{i, j}|$. The subset of activations in a Transformer block that are deemed outliers can be described by:
            \begin{equation}\label{eq:layer_activation}
                O = \{(i, j) \: | \: |h_{i, j}| > \tau \cdot \mu_h\},
            \end{equation}
            where $\tau$ is a predefined threshold.
            
            \item \textbf{Activation outliers in down-projection:} For the down-projection input $x_l^{down} \in \mathbb{R}^{B \times H}$ and the mean absolute value $\mu_{x^{down}} = \frac{1}{B \cdot H} \sum_{i, j} |x_{i, j}^{down}|$, the set of activation outliers in down-projection is given by:
            \begin{equation}\label{eq:down_proj_activation}
                O = \{(i, j) \: | \: |x_{i, j}^{down}| > \tau \cdot \mu_{x^{down}}\}
            \end{equation}
            
            \item \textbf{Weight outliers:} Considering a projection weight matrix $W \in \mathbb{R}^{O \times I}$, where $O$ and $I$ represent the output and input dimensions respectively, let $\mu_{w_{i}} = \frac{1}{I}\sum_{j}|w_{i, j}|$ denote the row-wise mean of the absolute values. The set of weight outliers can then be defined as:
            \begin{equation}\label{eq:weight_outliers}
                O = \{(i, j) \: | \: |w_{i, j}| > \tau \cdot \mu_{w_i}\}
            \end{equation}
            
            \item \textbf{Attention outliers:} For the attention matrix $A \in \mathbb{R}^{L \times L}$ with $L$ as the sequence length, $\hat{A_j} = \sum_{i = 1}^L A_{i, j}$ represents the total attention received by token $j$, and the mean attention value is $\mu_A = \frac{1}{L} \sum_{j} \hat{A}_j$. Attention outliers are categorized as:
            \begin{equation}\label{eq:attention_outliers}
                O = \{j \: | \: \hat{A}_j > \tau \cdot \mu_A\}
            \end{equation}
        \end{enumerate}

        These systematic outliers have a pronounced effect on model accuracy and complicate compression techniques like pruning and quantization \cite{an2025systematicoutlierslargelanguage}. 
        
        When pruning is considered, eliminating critical weight outliers can drastically harm performance, with accuracy reductions reaching up to 50\% for \llamaTwoSmall{} on few-shot tasks \cite{yu2024superweightlargelanguage}. On the other hand, maintaining these essential weights while removing less significant ones limits quality loss to under 2\% \cite{yu2024superweightlargelanguage}.

        Nevertheless, even the retention of important weights is not always sufficient to avoid degradation. Magnitude pruning \cite{han2016deepcompressioncompressingdeep}, which removes weights with the smallest absolute values (\ref{eq:magnitude}), still result in a 15\% drop in performance for \llamaTwoSmall{} \cite{wanda}.
        \begin{equation}\label{eq:magnitude}
            i_\xi = |W|,
        \end{equation}
        where $|W|$ is the absolute value of weight matrix.

        More sophisticated approaches, such as Wanda \cite{wanda} (\ref{eq:wanda}), which incorporates both weight magnitude and input activation, achieve improved sparsity and yield a 10\% perplexity enhancement over classic magnitude pruning for \llamaTwoSmall{}. However, a 5–15\% quality loss compared to the original dense architecture still persists.
        \begin{equation}\label{eq:wanda}
            i_T = |W| \cdot \|X_T\|_2,
        \end{equation}

        When quantization is considered, methods that prioritize significant weights have also been devised. Outlier-aware quantization \cite{yu2024superweightlargelanguage}, for example, surpasses SmoothQuant \cite{xiao2024smoothquantaccurateefficientposttraining} by 1\% on both \llamaTwoSmall{} and \llamaTwoBig{} \cite{yu2024superweightlargelanguage}. Furthermore, excluding these high-importance weights from quantization and keeping them in FP16 precision can boost model performance by as much as 67\% versus full quantization \cite{quantization_1}. Sensitivity-driven, non-uniform quantization, which allocates bits based on importance and preserves outliers in a sparse layout, can reduce quality loss to less than 2\% when moving from 16-bit to 4-bit formats \cite{kim2024squeezellmdenseandsparsequantization}.
    \subsubsection{Importance of layers}
        Recent research has delved deeply into the importance of different layers in Transformer architecture, uncovering nuanced insights about how these components function and contribute to overall model performance.

        Earlier investigations explored the robustness of Transformer layers in \llms{} by testing the effects of re-initializing or randomly resetting the parameters of specific layers while leaving the rest of the network untouched \cite{zhang2022layerscreatedequal}. These studies found that certain layers, termed as \textit{robust}, could be reset without noticeably degrading the model’s accuracy, implying that not every layer holds equal significance. Particularly in highly over-parameterized models trained with stochastic gradient descent, some layers naturally become less essential as the model restricts its own complexity \cite{zhang2022layerscreatedequal}.

        Moreover, analysis of individual Attention heads within the Transformer structure has shown that these heads are not interchangeable: they fulfill specialized roles during text generation \cite{voita2019analyzingmultiheadselfattentionspecialized}. For instance, some heads are crucial for recognizing syntactic patterns and dependency structures, highlighting their unique contributions. Interestingly, even after substantially reducing the number of Attention heads, the encoder’s performance remains largely unaffected, suggesting room for further model simplification \cite{voita2019analyzingmultiheadselfattentionspecialized}.

        Related findings indicate that the metric of perplexity, commonly used to assess model quality, is more strongly correlated with the total capacity of the network rather than the existence of self-attention modules alone \cite{liu2021payattentionmlps}. Thus, expanding model size or complexity may yield greater improvements than merely increasing self-attention layers.

        Building on this line of inquiry, Tyukin et al. \cite{tyukin2024attentionneeddontneed} demonstrated that pruning 40\% of the Multi-Head Attention layers leads to only minor performance drops, whereas removing a comparable proportion of MLP layers causes a much steeper decline up to 60\% in the \llamaTwoSmall{} model. This distinction highlights the foundational importance of MLP components in sustaining model efficacy. Further, He et al. \cite{he2024matterstransformersattentionneeded} identified that Multi-Head Attention layers contain a high degree of redundancy, meaning that many such layers can be omitted without significantly affecting results. This suggests that Transformer models could be made more efficient by reducing unnecessary Attention layers.
    
    \subsubsection{Importance of Transformer blocks}
        When considering larger components within \llm{}, one can evaluate the significance of entire Transformer block.

        Multiple techniques have been suggested for determining the importance of these blocks in \llms{}. Building on the principles introduced in the \wanda{} method \cite{wanda}, the \ows{} approach \cite{ows} performs selective pruning of Transformer blocks, guided by their outlier distributions. This method observes that most outliers are concentrated in the first and last blocks. Applying \ows{} enables a reduction in memory requirements by 12\% and delivers a slight, 3\% improvement in performance compared to \lora{}. Nonetheless, because \ows{} modifies a large number of parameters, it is not efficient in terms of computation time.

        According to UIDL \cite{uidl}, the layer-to-layer transformation of representations in a Transformer follows a residual iterative formula (\ref{eq:uidl_1}).
        \begin{equation}\label{eq:uidl_1}
            x^{(l + 1)} = x^{(l)} + f(x^{(l)}, \theta^{(l)}),
        \end{equation}
        where $x^{(l)}$ is the multi-dimensional input, $\theta^{(l)}$ is parameters vector for layer $l$, $f(x, \theta)$ describes the transformation of one Transformer block. 

        To ffectively remove the $n$ layers preceding layer  $l$ it is preferable if the input to the block about to be pruned closely matches its output (\ref{eq:uidl_2}). 
        \begin{equation}\label{eq:uidl_2}
            x^{(l)} = x^{(l - n)} + \varepsilon,
        \end{equation}
        where $\varepsilon$ is a constant indicating the degree to which one input may differ from another.

        Therefore, Transformer blocks that have the least impact on model predictions can be pinpointed by measuring the cosine similarity between their inputs and outputs (\ref{eq:uidl_3}).

        \begin{equation}\label{eq:uidl_3}
            i_T = argmin \frac{1}{\pi} arccos(\frac{X_T \cdot Y_T}{\|X_T\|_2 \cdot \|Y_T\|_2}),
        \end{equation}
        where $X_T$ and $Y_T$ are the input and output to the Transformer block, respectively, $\|\|_2$ is a Frobenius norm.

        This strategy achieves up to a 40\% reduction in the number of layers for the LLaMA2-70B model \cite{uidl} without sacrificing performance. However, \uidl{} involves extra inference passes on calibration samples, leading to higher computational costs. Additionally, a similar use of cosine similarity in ShortGPT \cite{men2024shortgptlayerslargelanguage} resulted in an 18\% speedup but at the expense of a 14\% drop in accuracy.

        Another methodology, \sleb{} \cite{sleb}, evaluates the relevance of each Transformer block by analyzing perplexity. The underlying assumption is that less essential blocks contribute minimally to the model’s predictions, so their removal causes only minor increases in perplexity (\ref{eq:sleb}).
    
        \begin{equation}\label{eq:sleb}
            i_T = exp\{-\frac{1}{S}\sum log p_\theta^n (X_T^{(S)}|X_{<T}^{(S)})\},
        \end{equation}
        where $X_T$ is the input to Transformer block, $S$ is a sequence length.

        Using \sleb{} yields a 27\% speed increase with an 8\% loss in quality for LLaMA2-7B compared to the original dense model \cite{sleb}. Like other methods, this approach also requires extra inference on a calibration set, extending overall processing time. Moreover, perplexity-based metrics are not suitable for scenarios where Transformer blocks need to be frozen dynamically, since they interfere with gradient computation.

        In the \shortllama{} method \cite{kim2024shortenedllamadepthpruning}, block importance is determined by measuring its effect on the loss function, essentially quantifying the change in loss when a block’s parameters are perturbed (\ref{eq:shortened_llama}).

        \begin{equation}\label{eq:shortened_llama}
            i_T = \mathit{L}(W_T + \Delta W_T) - \mathit{L}(W_T) = \nabla_W \mathit{L} \Delta W_T,
        \end{equation}
        where $W_T$ is a Transformer block weight, $L$ is a loss.

        This approach provides a 23\% speedup with a 7\% degradation in final quality for the LLaMA2-7B model \cite{kim2024shortenedllamadepthpruning}. However, as with \sleb{}, it suffers from additional computational overhead due to the need for calibration dataset inference. Also, loss-based metrics are unsuitable for on-the-fly removal of non-critical Transformer blocks.

    \subsubsection{Neural Architecture Search}
        Finally, it becomes possible to assess the value of entire architectures. Methods for this evaluation have been advanced under the area known as Neural Architecture Search (NAS). NAS refers to an automated process that designs neural network models, removing the need for manual architecture engineering. By employing NAS algorithms, it is feasible to automatically search for and identify the most suitable architectures tailored to specific applications.

        To achieve this, NAS typically employs a dedicated controller to gauge how well a candidate architecture might perform, often using a reduced dataset as a proxy. Nevertheless, this assessment technique, especially when relying on validation accuracy, can demand significant computational resources. For instance, the NASnet method \cite{nasnet} reportedly consumed over 2000 GPU days to find an optimal model structure.

        To mitigate such high computational requirements, zero-cost proxies have emerg-ed as a more efficient solution. These approaches use various heuristic measures to estimate the potential of candidate architectures without extensive training. For example, TF-TAS \cite{tf_tas} leverages synaptic diversity for the MHA component and synaptic saliency for the FFNN part. This method not only surpasses previous schemes, yielding a 4\% improvement in quality over hand-crafted models like ResNet and MobileNet-V3, but also drastically cuts down the computational resources needed, especially compared to approaches like AutoFormer \cite{autoformer_ti} and ViTAS \cite{vitas}.

        A detailed comparative analysis of 15 different zero-cost proxies applicable to Transformer-based architectures and \llms{} \cite{LPZero} demonstrated that their performance varies considerably. The effectiveness of these proxies was measured using Kendall and Spearman rank correlations, which range from -1 to 1, with higher scores reflecting better alignment between proxy-based rankings and those derived from fully trained networks. Despite their advantages, many proxies recorded modest correlation levels, with Kendall’s coefficients ranging from 0.01 up to 0.5, and Spearman’s from 0.1 to 0.5 \cite{LPZero}.

    \subsubsection{Summary}
        Therefore, various structures of \llms{}, including individual weights, layers, and Transformer blocks, as well as entire model architectures, can be evaluated in terms of their importance, specifically regarding their impact on the overall performance of the model. Such importance assessment methods are actively utilized in different optimization strategies for \llm{} deployment during inference, such as sparsification and quantization. Among the structural components considered, the Transformer block is particularly noteworthy. Existing research has already demonstrated the potential for compressing \llm{}, achieving up to a 30\% time and memory usage reduction by assessing the importance of Transformer blocks and retaining only the most significant ones. In contrast, approaches that focus on individual elements or layers have not achieved comparable resource savings, while methods focusing on entire architectures require substantial time and memory resources and utilized only within a Neural Architecture Search field. Hence, the analysis of Transformer block importance emerges as the most promising direction in this context. 

        However, while existing solutions for estimating the Transformer block importance successfully minimize computational costs, particularly  in the field of model inference, they often lead to a substantial quality degradation of an \llm{} performance. As a step towards addressing this challenge, the assessment of the importance of Transformer blocks can be further utilized within a fine-tuning field improving model quality, concurrently optimizing fine-tuning costs. 
        
        The next chapter presents a novel fine-tuning strategy, that dynamically allocates \lora{} adapters to essential Transformer blocks based on efficient importance analy-sis to further accelerate the fine-tuning process, while preserving the model quality.
